\section{Method}
\label{sec:method}

We formulate garment pattern layout as a mixed-type optimisation problem
and solve it with a multi-objective genetic algorithm.
The pipeline maps a genotype encoding seam positions, patch orientations,
texture phase offsets, and nesting parameters to a fabric layout,
which is then scored by three objectives: fabric efficiency,
texture alignment across seams, and garment fit quality.

% ---------------------------------------------------------------------------
\subsection{Problem Formulation}
\label{sec:formulation}

We define the problem as
\[
\Psi = \langle M, \mathcal{T}, \mathcal{L}, \mathbf{w}, \Omega \rangle,
\]
where $M = (V, E, F)$ is the 3D garment mesh,
$\mathcal{T}$ is a periodic fabric texture characterised by a lattice basis
$B = \{\mathbf{u}, \mathbf{v}\}$,
$\mathcal{L} = \{L_1, \dots, L_m\}$ is a set of admissible quad regions on
the mesh surface that constrain where seam keypoints may be placed,
$\mathbf{w} = \{w_s\}_{s \in \mathcal{S}}$ are static per-seam importance
weights ($w_s \in [0,1]$, set to zero for boundary seams),
and $\Omega = [0, W] \times [0, \infty)$ is the fabric roll domain of
width $W$.

% ---------------------------------------------------------------------------
\subsection{Texture Symmetry: Wallpaper Groups}
\label{sec:wallpaper}

The periodic structure of a fabric texture is fully characterised by its
wallpaper group $\mathcal{G}$, which encodes all translational, rotational,
and reflective symmetries of the pattern.
Accounting for $\mathcal{G}$ is essential for two reasons.
First, it constrains which grain-rotation pairs $(\rho_i, \rho_j)$ of adjacent
patches can produce a valid seam alignment: for a stripe pattern, rotating one
patch by 90° makes alignment impossible regardless of placement.
Second, patterns with glide-reflection symmetry (e.g.\ herringbone, chevron)
allow patches to align via a reflected copy of the texture, which must be
accounted for in the mismatch computation.

We support eight groups, summarised in Table~\ref{tab:wallpaper}.

\begin{table}[t]
\centering
\caption{Supported wallpaper groups with seam compatibility rules.}
\label{tab:wallpaper}
\begin{tabular}{llll}
\toprule
Group & Name & Pattern & Seam compatibility \\
\midrule
PM  & stripes          & horizontal stripes    & $\rho_i \equiv \rho_j \pmod{2}$ \\
PM  & diagonal stripes & stripes at 45°        & $\rho_i \equiv \rho_j \pmod{2}$ \\
PMM & grid             & check / plaid         & all pairs \\
P4  &                  & 4-fold rotation       & all pairs \\
P4M &                  & 4-fold + mirrors      & all pairs \\
PG  & herringbone      & glide reflection only & $\rho_i \equiv \rho_j \pmod{2}$ \\
PMG & chevron          & mirror + glide        & $\rho_i \equiv \rho_j \pmod{2}$ \\
PGG & brick bond       & two glide axes        & $\rho_i \equiv \rho_j \pmod{2}$ \\
\bottomrule
\end{tabular}
\end{table}

If $\mathcal{G}$ declares $(\rho_i, \rho_j)$ incompatible, the pair receives
a hard penalty in $f_2$.
For groups with glide-reflection symmetry (PG, PMG, PGG), each group defines
glide maps $\mathcal{G}_s = \{g_1, \dots\}$ acting on fractional phase
coordinates, and the per-point mismatch takes the minimum over the identity
and all glide images:
\[
\Delta_{\mathcal{G}}(\boldsymbol{\phi}_a, \boldsymbol{\phi}_b)
= \min\!\Bigl(
    \Delta(\boldsymbol{\phi}_a, \boldsymbol{\phi}_b),\;
    \min_{g \in \mathcal{G}_s} \Delta(\boldsymbol{\phi}_a, g(\boldsymbol{\phi}_b))
  \Bigr).
\]
For groups without glide symmetry, $\mathcal{G}_s = \emptyset$ and
$\Delta_{\mathcal{G}} = \Delta$.

% ---------------------------------------------------------------------------
\subsection{Genotype: Decision Variables}
\label{sec:genotype}

A candidate solution is encoded as
\[
\mathbf{x} =
\langle
\boldsymbol{\delta},\,
\boldsymbol{\rho},\,
\boldsymbol{\kappa},\,
\pi,\,
h,\,
\lambda_s
\rangle.
\]

\paragraph{Seam topology ($\boldsymbol{\delta}$).}
Each landmark region $L_j$ is a topological quad on the mesh surface defined
by four corner vertices.
The position of seam keypoint $k_j$ within $L_j$ is controlled by a 2D
parameter:
\[
k_j = \Gamma_j(\boldsymbol{\delta}_j), \quad
\boldsymbol{\delta}_j = (u_j, v_j) \in [0,1]^2,
\]
where $\Gamma_j$ performs bilinear interpolation within $L_j$ followed by a
snap to the nearest mesh vertex.
The full vector $\boldsymbol{\delta} \in [0,1]^{2m}$ concatenates all
landmark coordinates.
Seams $s \in \mathcal{S}$ are computed as geodesics between corresponding
keypoints.
Restricting keypoints to semantically defined regions prevents degenerate cuts
while still allowing the EA to explore meaningful seam placements.

\paragraph{Fabric orientation ($\boldsymbol{\rho}$).}
Cutting along seams yields $n$ 2D charts $P(\boldsymbol{\delta}) = \{p_1, \dots, p_n\}$.
Each chart is assigned a discrete grain orientation
\[
\rho_i \in \{0,1,2,3\}, \quad \theta_i = \rho_i \frac{\pi}{2},
\]
which determines how the chart is rotated on the fabric roll.
Compatibility with $\mathcal{G}$ is enforced at evaluation time.

\paragraph{Texture phase offsets ($\boldsymbol{\kappa}$).}
Because the fabric texture repeats periodically, placing a chart at different
positions that are integer multiples of the period apart produces
indistinguishable results.
We exploit this freedom by assigning each chart a discrete phase bin
\[
\kappa_i \in \{0,1,\dots,K-1\},
\]
where $K$ is the number of bins per period.
The bin determines a fractional offset $\kappa_i / K$ applied uniformly along
both lattice axes, effectively choosing which repeat of the texture the chart
is placed on.
Optimising $\boldsymbol{\kappa}$ allows the EA to align stripe or grid bands
across seams without moving the patches on the fabric.

\paragraph{Parameterization hyperparameter ($\lambda_s$).}
The seam straightness weight
\[
\lambda_s \in \{2^0, 2^1, \dots, 2^7\}
\]
controls the strength of the regularization term in the parameterization
energy that encourages seam boundaries to be straight lines in 2D.
It is applied uniformly across all seams.

\paragraph{Nesting variables ($\pi, h$).}
The permutation $\pi \in S_n$ defines the order in which charts are placed
on the fabric, and $h \in \mathcal{H}$ ($|\mathcal{H}|=3$) selects the
bottom-left placement heuristic variant.

% ---------------------------------------------------------------------------
\subsection{Phenotype: From Genotype to Layout}
\label{sec:phenotype}

\subsubsection{Seam Cutting and Parameterization.}
Given $\boldsymbol{\delta}$, the keypoints $\{k_j\}$ are computed via
bilinear interpolation and snapped to mesh vertices.
Geodesic paths between corresponding keypoints define the seam network,
along which the mesh is cut to yield 3D patches.
Each patch is then flattened to 2D by LSCM parameterization, with a
seam-straightness regularization weighted by $\lambda_s$.

\subsubsection{Stage 2: Global Phase Alignment.}
\label{sec:stage2}
After parameterization, patches connected by seams with $w_s > 0$ form
connected components.
Within each component a small translational correction $\mathbf{c}_i \in \mathbb{R}^2$
is found per patch (root fixed at $\mathbf{c}_i = \mathbf{0}$) by minimising
the total seam phase residual:
\[
C = \sum_{\substack{s \in \mathcal{S} \\ w_s > 0}}
    w_s \sum_{(a,b) \in s} \left\|\mathbf{r}_{ab}\right\|^2,
\]
where the signed wrapped residual per axis is
\[
r^d_{ab}
= \mathrm{wrap}\!\left(
    \mathrm{Phase}_i(\mathbf{x}_a + \mathbf{c}_i)^d
  - \mathrm{Phase}_j(\mathbf{x}_b + \mathbf{c}_j)^d
\right), \quad d \in \{u,v\},
\]
and $\mathrm{wrap}(\cdot)$ maps to $[-\tfrac{1}{2},\tfrac{1}{2})$.
This is solved with Levenberg--Marquardt using finite-difference Jacobians.
The corrections are baked into the patch vertices before nesting.

\subsubsection{Nesting.}
Each chart is rotated by $R(\theta_i)$ and placed on the fabric roll using
heuristic $h$ in order $\pi$.
Placements are snapped to the texture lattice so that the absolute position
on the roll does not introduce an uncontrolled phase offset.
Charts must satisfy non-overlap and containment constraints:
\begin{align}
(\phi_i(p_i) + \mathbf{t}_i) \cap (\phi_j(p_j) + \mathbf{t}_j) &= \emptyset, \\
(\phi_i(p_i) + \mathbf{t}_i) &\subset \Omega.
\end{align}

% ---------------------------------------------------------------------------
\subsection{Objective Functions}
\label{sec:objectives}

The fitness vector is
\[
\mathbf{F}(\mathbf{x}) = [w_1 f_1,\; w_2 f_2,\; w_3 f_3]^T,
\]
where $w_1, w_2, w_3 \geq 0$ are designer-specified weights and all three
components are minimised.

\paragraph{Fabric efficiency ($f_1$).}
\[
f_1(\mathbf{x}) = \frac{1}{W} \max_{p \in \text{Layout}} y_{\max}(p).
\]
Normalising by $W$ makes $f_1$ dimensionless and comparable to $f_2$
regardless of the physical scale of the fabric.

\paragraph{Texture alignment ($f_2$).}
The 2D lattice phase of a fabric point $\mathbf{x}$ is
\[
\mathrm{Phase}(\mathbf{x}) = \mathrm{frac}([B]^{-1}\mathbf{x}) \in [0,1)^2,
\]
where $[B] = [\mathbf{u} \mid \mathbf{v}]$ and $\mathrm{frac}(\cdot)$ is
applied component-wise.
For chart $p_i$, the kappa offset shifts the phase uniformly:
\[
\mathrm{Phase}_i(\mathbf{x}) =
\mathrm{frac}\!\left(\mathrm{Phase}(\mathbf{x}) + \frac{\kappa_i}{K}\mathbf{1}\right),
\quad \mathbf{1} = (1,1)^T.
\]
For stripe groups, the pattern is translation-invariant along the stripe axis,
so that phase component contributes zero to any difference.
The per-point mismatch between two phases is
\[
\Delta(\boldsymbol{\phi}_a, \boldsymbol{\phi}_b)
= \tfrac{1}{2}\bigl(
    \delta^u(\boldsymbol{\phi}_a, \boldsymbol{\phi}_b)
  + \delta^v(\boldsymbol{\phi}_a, \boldsymbol{\phi}_b)
\bigr),
\quad
\delta^d = \min\!\left(|\phi_a^d - \phi_b^d|,\; 1 - |\phi_a^d - \phi_b^d|\right).
\]
The seam mismatch integrates this over each seam arc-length:
\[
\mathrm{Mismatch}(s)
= \frac{1}{L_s} \int_0^{L_s}
  \Delta_{\mathcal{G}}\!\bigl(
    \mathrm{Phase}_i(\mathbf{x}_i(l)),\,
    \mathrm{Phase}_j(\mathbf{x}_j(l))
  \bigr)\, dl,
\]
and $f_2$ accumulates this over all seams weighted by importance:
\[
f_2(\mathbf{x}) = \sum_{s \in \mathcal{S}} w_s \cdot \mathrm{Mismatch}(s).
\]

\paragraph{Fit distortion ($f_3$).}
\[
f_3(\mathbf{x}) = \frac{1}{n} \sum_{i=1}^n \overline{\sigma}_i^{\,\mathrm{weft}},
\]
where $\overline{\sigma}_i^{\,\mathrm{weft}}$ is the mean local triangle
stretch deviation along the weft direction between corresponding 2D and 3D
triangles of chart $p_i$.

% ---------------------------------------------------------------------------
\subsection{Evolutionary Algorithm}
\label{sec:ea}

The mixed structure of the genotype — continuous ($\boldsymbol{\delta}$),
discrete ($\boldsymbol{\rho}, \boldsymbol{\kappa}, \lambda_s, h$),
and combinatorial ($\pi$) — motivates a genetic algorithm with
variable-specific operators rather than a standard real-valued EA.

\paragraph{Selection.}
Tournament selection with tournament size $k=3$ is used.
The primary criterion is Pareto dominance: a candidate $\mathbf{x}$ dominates
$\mathbf{y}$ if $\mathbf{F}(\mathbf{x}) \leq \mathbf{F}(\mathbf{y})$
component-wise with at least one strict inequality.
When neither candidate dominates the other, the fallback criterion is the
scalar sum $\sum_i w_i f_i$, which provides a consistent total order
without discarding multi-objective information.

\paragraph{Crossover.}
Applied with probability $p_c$.
Each variable type uses a dedicated operator:
\begin{itemize}
    \item $\boldsymbol{\delta}, \boldsymbol{\rho}, \boldsymbol{\kappa}$:
          uniform gene-wise swap — each gene is exchanged between parents
          independently with probability $\tfrac{1}{2}$.
    \item $\pi$: whole-permutation inheritance — one parent's permutation is
          copied in full (chosen by coin flip) to preserve permutation validity.
    \item $h, \lambda_s$: coin flip between the two parent values.
\end{itemize}

\paragraph{Mutation.}
Applied with probability $p_m$ after crossover.
Each variable type uses a dedicated operator:
\begin{itemize}
    \item $\boldsymbol{\delta}$: Gaussian perturbation
          $\delta_j \leftarrow \mathrm{clip}(\delta_j + \mathcal{N}(0, \sigma_\delta),\,
          \delta_{\min}, \delta_{\max})$,
          where $[\delta_{\min}, \delta_{\max}] \subset [0,1]$ avoids
          geometrically degenerate mesh regions.
    \item $\boldsymbol{\kappa}$: each bin resampled uniformly from
          $\{0,\dots,K{-}1\}$ with per-gene probability $p_\kappa$.
    \item $\boldsymbol{\rho}$: each value resampled from $\{0,1,2,3\}$
          with per-gene probability $p_\rho$.
    \item $\pi$: swap of two randomly chosen positions with probability
          $p_\pi$, preserving permutation validity.
    \item $h$: resampled uniformly from $\mathcal{H}$ with probability $p_h$.
    \item $\lambda_s$: resampled uniformly from $\{2^0,\dots,2^7\}$
          with probability $p_{\lambda}$.
\end{itemize}

\paragraph{Elitism.}
The single best individual (by fitness sum) is carried over unchanged to the
next generation, ensuring monotone improvement of the best solution found.
